% ============================================================
% PART 4: UC23 - UC30
% ============================================================

\subsection*{Use Case 23: Ghi nhận số tiền tip}
\begin{longtable}{|P{4cm}|P{11cm}|}
\hline
\textbf{Thuộc tính} & \textbf{Chi tiết} \\ \hline \endfirsthead
\hline \textbf{Thuộc tính} & \textbf{Chi tiết} \\ \hline \endhead
Tên Use Case & UC23. Ghi nhận số tiền tip của khách hàng vào hệ thống thanh toán \\ \hline
Actor chính & Thu ngân (Cashiers) \\ \hline
Mục đích & Lưu vết số tiền khách bo (tip) thêm ngoài hóa đơn để báo cáo doanh thu và phân chia cho nhân viên vào cuối ngày. \\ \hline
Tiền điều kiện & Đơn hàng đang được thanh toán bằng Thẻ tín dụng hoặc khách đưa dư tiền mặt và bảo giữ lại làm tip. \\ \hline
Hậu điều kiện & Số tiền Tip được ghi nhận riêng biệt so với Doanh thu bán hàng, phục vụ cho báo cáo tài chính. \\ \hline
Luồng sự kiện chính &
1. Trước khi chốt giao dịch thẻ, hệ thống hiển thị màn hình yêu cầu nhập số tiền Tip.\newline
2. Khách hàng chọn \% tip hoặc nhập số tiền trên màn hình phụ (Customer Display).\newline
3. Thu ngân (hoặc khách) nhập con số đó vào ô ``Tiền Tip'' trên POS.\newline
4. Hệ thống cộng dồn tiền món ăn và tiền tip để ra tổng số tiền thực tế trừ vào thẻ.\newline
5. Hệ thống ghi nhận giao dịch: Tách rõ phần ``Doanh thu món'' và ``Quỹ Tip''. \\ \hline
Luồng ngoại lệ &
\textbf{3a. Nhập nhầm tiền Tip quá lớn:} Nếu Tip lớn hơn 50\% giá trị hóa đơn, hệ thống bật popup cảnh báo ``Số tiền tip bất thường. Xác nhận?''. Thu ngân phải kiểm tra lại với khách trước khi tiếp tục. \\ \hline
\end{longtable}

\subsection*{Use Case 24: Thêm, sửa, xóa món ăn trong thực đơn}
\begin{longtable}{|P{4cm}|P{11cm}|}
\hline
\textbf{Thuộc tính} & \textbf{Chi tiết} \\ \hline \endfirsthead
\hline \textbf{Thuộc tính} & \textbf{Chi tiết} \\ \hline \endhead
Tên Use Case & UC24. Thêm, sửa, xóa các món ăn trong thực đơn và cấu hình giá bán \\ \hline
Actor chính & Quản lý (Managers) \\ \hline
Mục đích & Cập nhật bộ thực đơn cốt lõi của nhà hàng (thêm món mới theo mùa, đổi giá) để POS và KDS đồng bộ ngay lập tức. \\ \hline
Tiền điều kiện & Quản lý đã đăng nhập vào hệ thống Admin Web với quyền ``Admin'' hoặc ``Manager''. \\ \hline
Hậu điều kiện & Cơ sở dữ liệu thực đơn được cập nhật; mọi thiết bị POS/App tự động tải về thực đơn mới nhất. \\ \hline
Luồng sự kiện chính &
1. Quản lý truy cập mục ``Quản lý Thực đơn'' (Menu Management).\newline
2. Quản lý nhấn ``Thêm món mới'' hoặc chọn một món hiện có để ``Sửa''.\newline
3. Hệ thống hiển thị biểu mẫu chi tiết: Tên món, Hình ảnh, Mô tả, Giá bán, Phân loại, Trạm bếp phụ trách.\newline
4. Quản lý nhập thông tin và nhấn ``Lưu thay đổi''.\newline
5. Hệ thống lưu vào Database và phát broadcast event báo cho tất cả thiết bị POS tải lại cache (Menu sync). \\ \hline
Luồng ngoại lệ &
\textbf{4a. Xóa món đang nằm trong đơn hàng chưa phục vụ xong:} Hệ thống chặn lệnh Xóa và báo lỗi, yêu cầu chuyển sang trạng thái ``Ẩn/Hết hàng'' tạm thời. \\ \hline
\end{longtable}

\subsection*{Use Case 25: Đánh dấu món ăn ``hết hàng'' theo thời gian thực}
\begin{longtable}{|P{4cm}|P{11cm}|}
\hline
\textbf{Thuộc tính} & \textbf{Chi tiết} \\ \hline \endfirsthead
\hline \textbf{Thuộc tính} & \textbf{Chi tiết} \\ \hline \endhead
Tên Use Case & UC25. Đánh dấu món ăn ``hết hàng'' trên thực đơn theo thời gian thực \\ \hline
Actor chính & Quản lý (Managers) / Đầu bếp trưởng (Head Chef) \\ \hline
Mục đích & Ngăn chặn nhân viên phục vụ tiếp tục order món ăn đã hết nguyên liệu dưới bếp, tránh việc phải ra bàn xin lỗi khách. \\ \hline
Tiền điều kiện & Người dùng có quyền ``Cập nhật trạng thái thực đơn''. \\ \hline
Hậu điều kiện & Món ăn bị làm mờ (gray out) trên tất cả các màn hình gọi món, nhân viên không thể bấm chọn. \\ \hline
Luồng sự kiện chính &
1. Đầu bếp báo hết nguyên liệu cho Quản lý, hoặc Quản lý tự nhận thấy qua cảnh báo tồn kho.\newline
2. Quản lý mở POS hoặc Admin Web, vào chức năng ``Trạng thái Menu 86'' (Menu 86 / Sold out).\newline
3. Quản lý tìm món ăn và bật công tắc từ ``Còn hàng'' (In stock) sang ``Hết hàng'' (Sold out).\newline
4. Hệ thống ngay lập tức đẩy bản cập nhật xuống tất cả các iPad/POS đang mở.\newline
5. Món ăn tự động hiển thị mờ đi kèm chữ ``Hết hàng'' không thể bấm chọn. \\ \hline
Luồng ngoại lệ &
\textbf{3a. Hết hàng có định mức (Chỉ còn 5 phần):} Quản lý chọn ``Giới hạn số lượng'' và nhập số ``5''. Hệ thống đếm lùi trên POS mỗi khi nhân viên order món đó, về 0 sẽ tự động khóa lại. \\ \hline
\end{longtable}

\section*{1.13. Đặc tả Use Case Scenario (Phần 6: UC 26--30)}

\subsection*{Use Case 26: Theo dõi cảnh báo hàng tồn kho dưới ngưỡng an toàn}
\begin{longtable}{|P{4cm}|P{11cm}|}
\hline
\textbf{Thuộc tính} & \textbf{Chi tiết} \\ \hline \endfirsthead
\hline \textbf{Thuộc tính} & \textbf{Chi tiết} \\ \hline \endhead
Tên Use Case & UC26. Theo dõi cảnh báo hàng tồn kho khi mức nguyên liệu xuống dưới ngưỡng an toàn \\ \hline
Actor chính & Quản lý (Managers) \\ \hline
Mục đích & Đảm bảo nhà hàng luôn có đủ nguyên liệu bằng cách nhắc nhở Quản lý đặt hàng kịp thời trước khi cạn kiệt. \\ \hline
Tiền điều kiện & Module Quản lý Tồn kho đã được thiết lập mức ``Ngưỡng an toàn'' (Reorder level) cho từng loại nguyên liệu. \\ \hline
Hậu điều kiện & Quản lý nhận được thông tin để lên kế hoạch nhập hàng mới. \\ \hline
Luồng sự kiện chính &
1. Hệ thống liên tục ghi nhận lượng nguyên liệu bị trừ đi sau mỗi lần bán món hoặc trừ thủ công từ bếp.\newline
2. Ngay khi số lượng nguyên liệu giảm xuống dưới ngưỡng an toàn, hệ thống tạo cảnh báo tự động.\newline
3. Hệ thống hiển thị cảnh báo trên màn hình Dashboard của Admin Web.\newline
4. (Tuỳ chọn) Hệ thống gửi email hoặc thông báo qua App quản lý nội dung.\newline
5. Quản lý xem cảnh báo, nhấn vào để xem lịch sử tiêu thụ và dự đoán thời gian cạn kiệt. \\ \hline
Luồng ngoại lệ &
\textbf{1a. Lỗi sai lệch định lượng:} Nếu tồn kho thực tế khác với hệ thống (do hao hụt không khai báo), cảnh báo có thể hiển thị muộn. Quản lý phải dùng chức năng ``Kiểm kê kho'' (Stock-take) để điều chỉnh lại. \\ \hline
\end{longtable}

\subsection*{Use Case 27: Xem báo cáo phân tích doanh thu}
\begin{longtable}{|P{4cm}|P{11cm}|}
\hline
\textbf{Thuộc tính} & \textbf{Chi tiết} \\ \hline \endfirsthead
\hline \textbf{Thuộc tính} & \textbf{Chi tiết} \\ \hline \endhead
Tên Use Case & UC27. Xem báo cáo phân tích doanh thu theo giờ cao điểm và món bán chạy \\ \hline
Actor chính & Quản lý (Managers) \\ \hline
Mục đích & Cung cấp cái nhìn tổng quan về hiệu quả kinh doanh, giúp Quản lý đưa ra quyết định thay đổi menu hoặc tăng cường nhân sự. \\ \hline
Tiền điều kiện & Quản lý đã đăng nhập; có đủ dữ liệu giao dịch từ module Thanh toán trong khoảng thời gian muốn xem. \\ \hline
Hậu điều kiện & Báo cáo trực quan (biểu đồ) được hiển thị hoặc xuất ra file thành công. \\ \hline
Luồng sự kiện chính &
1. Quản lý truy cập vào module ``Báo cáo \& Phân tích'' (Analytics \& Reports).\newline
2. Chọn loại báo cáo ``Doanh thu'' và thiết lập bộ lọc (Hôm nay, Tuần này, Tháng này).\newline
3. Hệ thống truy vấn dữ liệu từ cơ sở dữ liệu báo cáo (Reporting DB) để không ảnh hưởng đến DB giao dịch.\newline
4. Hệ thống vẽ biểu đồ doanh thu theo khung giờ và biểu đồ Top 5 món bán chạy nhất.\newline
5. Quản lý xem xét biểu đồ, có thể di chuột vào để xem con số chi tiết. \\ \hline
Luồng ngoại lệ &
\textbf{3a. Dữ liệu quá lớn gây tải chậm:} Hệ thống hiển thị thanh tiến trình (Loading bar) và thông báo ``Đang tổng hợp dữ liệu, vui lòng chờ...''. \\ \hline
\end{longtable}

\subsection*{Use Case 28: Phân tích hiệu suất nhân viên và độ trễ phục vụ}
\begin{longtable}{|P{4cm}|P{11cm}|}
\hline
\textbf{Thuộc tính} & \textbf{Chi tiết} \\ \hline \endfirsthead
\hline \textbf{Thuộc tính} & \textbf{Chi tiết} \\ \hline \endhead
Tên Use Case & UC28. Phân tích đánh giá hiệu suất nhân viên và độ trễ phục vụ tại bếp \\ \hline
Actor chính & Quản lý (Managers) \\ \hline
Mục đích & Phát hiện các nút thắt cổ chai (bottlenecks) trong vận hành, đánh giá đầu bếp nào làm việc nhanh nhất và nhân viên xử lý đơn tốt nhất. \\ \hline
Tiền điều kiện & Hệ thống lưu trữ đầy đủ tem thời gian (timestamps) của mọi hành động trên POS và KDS. \\ \hline
Hậu điều kiện & Báo cáo hiệu suất vận hành (Operational Analytics) được xuất ra chi tiết. \\ \hline
Luồng sự kiện chính &
1. Quản lý chọn tab ``Hiệu suất Vận hành'' trong module Báo cáo.\newline
2. Hệ thống tổng hợp các chỉ số: Thời gian trung bình từ ``Chờ chế biến'' đến ``Sẵn sàng'' (Prep Time) và từ ``Sẵn sàng'' đến ``Đã phục vụ'' (Service Time).\newline
3. Hệ thống nhóm dữ liệu theo từng Trạm bếp hoặc từng Nhân viên (Staff ID).\newline
4. Hệ thống hiển thị bảng xếp hạng, tô đỏ các trạm có tỷ lệ trễ SLA cao nhất.\newline
5. Quản lý dùng thông tin này để đào tạo lại hoặc sắp xếp lại nhân sự trong ca làm. \\ \hline
Luồng ngoại lệ &
\textbf{2a. Dữ liệu nhiễu (Nhân viên quên bấm nút hoàn thành):} Hệ thống tự động loại bỏ các giá trị ngoại lai (outliers) khỏi công thức tính trung bình để báo cáo không bị sai lệch. \\ \hline
\end{longtable}

\subsection*{Use Case 29: Phân quyền và cấp tài khoản truy cập}
\begin{longtable}{|P{4cm}|P{11cm}|}
\hline
\textbf{Thuộc tính} & \textbf{Chi tiết} \\ \hline \endfirsthead
\hline \textbf{Thuộc tính} & \textbf{Chi tiết} \\ \hline \endhead
Tên Use Case & UC29. Phân quyền và cấp tài khoản truy cập cho từng vai trò nhân sự \\ \hline
Actor chính & Quản lý (Managers) / Quản trị hệ thống (Admin) \\ \hline
Mục đích & Đảm bảo tính bảo mật của hệ thống bằng cách kiểm soát quyền truy cập dựa trên vai trò (Role-Based Access Control -- RBAC). \\ \hline
Tiền điều kiện & Quản lý đăng nhập bằng tài khoản có đặc quyền cao nhất (Super Admin / Manager). \\ \hline
Hậu điều kiện & Tài khoản mới được tạo hoặc quyền của tài khoản cũ được thay đổi, có hiệu lực ngay trong lần đăng nhập tiếp theo. \\ \hline
Luồng sự kiện chính &
1. Quản lý truy cập chức năng ``Quản lý Nhân sự / Tài khoản''.\newline
2. Quản lý nhấn ``Tạo tài khoản mới'' và nhập thông tin (Tên, Mã NV, Mật khẩu tạm).\newline
3. Quản lý gán vai trò (Role) cho tài khoản, ví dụ chọn ``Thu ngân''.\newline
4. Hệ thống tự động ánh xạ vai trò với các nhóm quyền cụ thể.\newline
5. Quản lý nhấn ``Lưu''.\newline
6. Hệ thống tạo tài khoản trong cơ sở dữ liệu và nhân viên có thể dùng ngay lập tức. \\ \hline
Luồng ngoại lệ &
\textbf{3a. Trùng lặp ID Nhân viên:} Hệ thống cảnh báo ``Mã Nhân viên này đã tồn tại trong hệ thống. Vui lòng nhập mã khác''. Khóa nút Lưu cho đến khi nhập đúng. \\ \hline
\end{longtable}

\subsection*{Use Case 30: Truy xuất nhật ký kiểm toán (Audit Logs)}
\begin{longtable}{|P{4cm}|P{11cm}|}
\hline
\textbf{Thuộc tính} & \textbf{Chi tiết} \\ \hline \endfirsthead
\hline \textbf{Thuộc tính} & \textbf{Chi tiết} \\ \hline \endhead
Tên Use Case & UC30. Truy xuất nhật ký kiểm toán (Audit Logs) để rà soát các thao tác hủy đơn và hoàn tiền \\ \hline
Actor chính & Quản lý (Managers) / Kế toán (Accountant) \\ \hline
Mục đích & Phát hiện và ngăn ngừa gian lận, đảm bảo tính minh bạch (ví dụ: nhân viên cố tình xóa đơn hàng đã thu tiền khách). \\ \hline
Tiền điều kiện & Người dùng có quyền truy cập module ``Nhật ký Hệ thống'' (Audit Logs). Hệ thống đã bật chức năng ghi log cho mọi thao tác nhạy cảm. \\ \hline
Hậu điều kiện & Quản lý xem được lịch sử chi tiết ``Ai đã làm gì, vào lúc nào'' để đối soát. \\ \hline
Luồng sự kiện chính &
1. Quản lý truy cập chức năng ``Nhật ký Kiểm toán''.\newline
2. Quản lý thiết lập bộ lọc: Hành động = ``Hủy đơn'' (Void) hoặc ``Hoàn tiền'' (Refund); Khung thời gian = ``Hôm qua''.\newline
3. Hệ thống trả về danh sách các sự kiện khớp với bộ lọc.\newline
4. Quản lý nhấn vào một dòng bất kỳ để xem chi tiết.\newline
5. Hệ thống hiển thị: Người thực hiện, Thời gian chính xác, Lý do hủy, ID Đơn hàng bị ảnh hưởng, và Quản lý nào đã duyệt. \\ \hline
Luồng ngoại lệ &
\textbf{3a. Hành vi đáng ngờ số lượng lớn:} Nếu một nhân viên có hơn 3 lượt hủy đơn trong cùng 1 ca làm, hệ thống bôi đỏ tên nhân viên đó trên dòng log để Quản lý đặc biệt chú ý kiểm tra. \\ \hline
\end{longtable}
